\documentclass{csc2059practical}

\setmodulecode{CSC2059}
\setmoduletitle{Data Structures and Algorithms}
\setpracticalnumber{2}
\setpracticaltitle{Introduction to C++ Functions, Compilation and Testing}
\setpracticalsubtitle{C++ functions, automated public tests, and the CSC2059 development workflow}
\setpracticaldate{}




% Phase 2 visual aids: drawn directly in LaTeX, not external images.
% Repository structure uses dirtree for a true-to-life directory tree.
% Workflow and compilation model use TikZ.
\usepackage{dirtree}
\usepackage{tabularx}
\usetikzlibrary{positioning,fit,backgrounds}

\newcommand{\classroomlink}
{
  https://example.com
}

\newcommand{\CSCWorkflowDiagram}{%
\begin{figure}[htbp]
\centering
\begin{tikzpicture}[
  stage/.style={circle, draw=##1!75!black, fill=##1!8, very thick,
    minimum size=21mm, align=center, font=\sffamily\bfseries\scriptsize},
  flow/.style={-{Latex[length=3mm]}, very thick, draw=CSCBlue!70},
  captionnote/.style={font=\sffamily\small, align=center, text=black!80}
]
\node[stage=CSCBlue] (read) at (90:28mm) {Read\\task};
\node[stage=CSCGreen] (edit) at (30:28mm) {Edit\\code};
\node[stage=CSCAmber] (test) at (-30:28mm) {Run\\tests};
\node[stage=CSCRed] (fix) at (-90:28mm) {Interpret\\errors};
\node[stage=CSCPurple] (commit) at (-150:28mm) {Commit};
\node[stage=CSCBlue] (push) at (150:28mm) {Push};

\draw[flow] (read) -- (edit);
\draw[flow] (edit) -- (test);
\draw[flow] (test) -- (fix);
\draw[flow] (fix) -- (commit);
\draw[flow] (commit) -- (push);
\draw[flow] (push) -- (read);

\node[captionnote] at (0,0) {repeat\\in small\\steps};
\end{tikzpicture}
\caption{The CSC2059 edit--test--commit workflow.}
\label{fig:workflow-cycle}
\end{figure}%
}

\newcommand{\CSCRepositoryTree}{%
\begin{figure}[htbp]
\centering
\begin{tcolorbox}[
  enhanced,
  colback=CSCAmber!4,
  colframe=CSCAmber!65!black,
  boxrule=0.6pt,
  arc=2mm,
  left=4mm,right=4mm,top=3mm,bottom=3mm,
  title={Repository structure},
  fonttitle=\sffamily\bfseries,
  coltitle=black,
  colbacktitle=CSCAmber!16,
]
{\ttfamily\small
\dirtree{%
.1 practical-02/.
.2 README.md.
.2 Makefile.
.2 myFunctions.h.
.2 myFunctions.cpp.
.2 main.cpp.
.2 tests/.
.3 public\_tests.cpp.
.2 .github/.
.3 workflows/.
.2 .classroom/.
}}
\medskip
\begin{tabularx}{\linewidth}{@{}>{\ttfamily\small}l X@{}}
README.md & Practical overview and reminders.\\
Makefile & Shortcuts for compiling, running, and testing the project.\\
myFunctions.h & Function declarations: what functions exist and what they return.\\
myFunctions.cpp & Function implementations: the code you will mostly edit.\\
main.cpp & Small program used for manual experiments and interaction.\\
tests/public\_tests.cpp & Automated public tests run using \texttt{make test}.\\
.github/, .classroom/ & Classroom50/GitHub configuration. Do not modify unless instructed.\\
\end{tabularx}
\end{tcolorbox}
\caption{The main files and folders in the Practical 2 repository.}
\label{fig:repository-tree}
\end{figure}%
}




\newcommand{\CSCBuildPipelineRun}{%
\begin{figure}[p]
\centering
\begin{tcolorbox}[
  enhanced,
  colback=CSCBlue!2,
  colframe=CSCBlue!55!black,
  boxrule=0.6pt,
  arc=2mm,
  left=5mm,right=5mm,top=5mm,bottom=5mm,
  title={What happens when you run \texttt{make run}},
  fonttitle=\sffamily\bfseries,
  coltitle=black,
  colbacktitle=CSCBlue!10,
]
\centering
\begin{tikzpicture}[
  sourceblue/.style={draw=CSCBlue!70!black, fill=CSCBlue!6, rounded corners=2mm, very thick,
    minimum width=39mm, minimum height=15mm, align=center, font=\sffamily\small},
  sourcegreen/.style={draw=CSCGreen!70!black, fill=CSCGreen!6, rounded corners=2mm, very thick,
    minimum width=39mm, minimum height=15mm, align=center, font=\sffamily\small},
  sourceamber/.style={draw=CSCAmber!70!black, fill=CSCAmber!8, rounded corners=2mm, very thick,
    minimum width=39mm, minimum height=15mm, align=center, font=\sffamily\small},
  process/.style={draw=CSCBlue!75!black, fill=CSCLightBlue, rounded corners=2mm, very thick,
    minimum width=62mm, minimum height=15mm, align=center, font=\sffamily\small\bfseries},
  output/.style={draw=CSCGreen!75!black, fill=CSCGreen!7, rounded corners=2mm, very thick,
    minimum width=62mm, minimum height=15mm, align=center, font=\sffamily\small\bfseries},
  note/.style={font=\sffamily\scriptsize, align=center, text=black!75, text width=118mm},
  arrow/.style={-{Latex[length=3mm]}, very thick, draw=CSCBlue!75!black}
]

\node[sourceblue] (header) at (-48mm,0) {\texttt{myFunctions.h}\\\footnotesize function declarations};
\node[sourcegreen] (library) at (0,0) {\texttt{myFunctions.cpp}\\\footnotesize function implementations};
\node[sourceamber] (main) at (48mm,0) {\texttt{main.cpp}\\\footnotesize manual program};

\node[process] (make) at (0,-24mm) {\texttt{make run}\\\footnotesize asks the Makefile to build the program};
\node[process] (compile) at (0,-48mm) {compile source files\\\footnotesize translate C++ source into object code};
\node[process] (link) at (0,-72mm) {link compiled parts\\\footnotesize combine the program and library code};
\node[output] (exe) at (0,-96mm) {run the program\\\footnotesize interactive input/output from \texttt{main.cpp}};
\node[note] (note) at (0,-122mm) {Use \texttt{make run} when you want to try the example program manually. It builds code from \texttt{main.cpp} together with the functions implemented in \texttt{myFunctions.cpp}.};

\draw[arrow] (header.south) |- (make.west);
\draw[arrow] (library.south) -- (make.north);
\draw[arrow] (main.south) |- (make.east);
\draw[arrow] (make) -- (compile);
\draw[arrow] (compile) -- (link);
\draw[arrow] (link) -- (exe);
\end{tikzpicture}
\end{tcolorbox}
\caption{The \texttt{make run} path builds the student program from \texttt{main.cpp} and the function library.}
\label{fig:make-run-model}
\end{figure}%
}

\newcommand{\CSCBuildPipelineTest}{%
\begin{figure}[p]
\centering
\begin{tcolorbox}[
  enhanced,
  colback=CSCGreen!2,
  colframe=CSCGreen!55!black,
  boxrule=0.6pt,
  arc=2mm,
  left=5mm,right=5mm,top=5mm,bottom=5mm,
  title={What happens when you run \texttt{make test}},
  fonttitle=\sffamily\bfseries,
  coltitle=black,
  colbacktitle=CSCGreen!10,
]
\centering
\begin{tikzpicture}[
  sourceblue/.style={draw=CSCBlue!70!black, fill=CSCBlue!6, rounded corners=2mm, very thick,
    minimum width=39mm, minimum height=15mm, align=center, font=\sffamily\small},
  sourcegreen/.style={draw=CSCGreen!70!black, fill=CSCGreen!6, rounded corners=2mm, very thick,
    minimum width=39mm, minimum height=15mm, align=center, font=\sffamily\small},
  sourcepurple/.style={draw=CSCPurple!70!black, fill=CSCPurple!6, rounded corners=2mm, very thick,
    minimum width=39mm, minimum height=15mm, align=center, font=\sffamily\small},
  process/.style={draw=CSCBlue!75!black, fill=CSCLightBlue, rounded corners=2mm, very thick,
    minimum width=62mm, minimum height=15mm, align=center, font=\sffamily\small\bfseries},
  output/.style={draw=CSCPurple!75!black, fill=CSCPurple!7, rounded corners=2mm, very thick,
    minimum width=62mm, minimum height=15mm, align=center, font=\sffamily\small\bfseries},
  note/.style={font=\sffamily\scriptsize, align=center, text=black!75, text width=118mm},
  arrow/.style={-{Latex[length=3mm]}, very thick, draw=CSCBlue!75!black}
]

\node[sourceblue] (header) at (-48mm,0) {\texttt{myFunctions.h}\\\footnotesize function declarations};
\node[sourcegreen] (library) at (0,0) {\texttt{myFunctions.cpp}\\\footnotesize function implementations};
\node[sourcepurple] (tests) at (48mm,0) {\texttt{tests/public\_tests.cpp}\\\footnotesize automated checks};

\node[process] (make) at (0,-24mm) {\texttt{make test}\\\footnotesize asks the Makefile to build the tests};
\node[process] (compile) at (0,-48mm) {compile source files\\\footnotesize translate your library and test code};
\node[process] (link) at (0,-72mm) {link test executable\\\footnotesize combine the tests with your functions};
\node[output] (result) at (0,-96mm) {run public tests\\\footnotesize report pass/fail feedback};
\node[note] (note) at (0,-122mm) {Use \texttt{make test} when you want automated feedback. The tests call your functions directly and compare the observed behaviour with the expected behaviour.};

\draw[arrow] (header.south) |- (make.west);
\draw[arrow] (library.south) -- (make.north);
\draw[arrow] (tests.south) |- (make.east);
\draw[arrow] (make) -- (compile);
\draw[arrow] (compile) -- (link);
\draw[arrow] (link) -- (result);
\end{tikzpicture}
\end{tcolorbox}
\caption{The \texttt{make test} path builds and runs the automated public tests against the same function library.}
\label{fig:make-test-model}
\end{figure}%
}

\newcommand{\CSCCompilationModel}{%
\clearpage
\CSCBuildPipelineRun
\CSCBuildPipelineTest
\clearpage
}

\begin{document}

\maketitle
\tableofcontents
\clearpage

\section{Introduction}

Welcome to your first C++ programming practical in CSC2059.

In the previous practical, you learned how to use Git, GitHub, and Visual Studio Code to manage your work. In this practical, you will begin writing and testing C++ code using the same development workflow that will be used throughout the module.

Rather than creating an entire project from scratch, you will work from a provided code scaffold. This reflects how software development often happens in practice: you are frequently asked to understand, modify, extend, and improve an existing code base rather than starting with an empty folder.

You will learn how C++ programs are organised across multiple files, how functions are declared and implemented, and how automated tests can help you verify that your code behaves correctly.

Throughout this practical you should work in small steps:

\begin{itemize}
  \item read the task;
  \item make a change;
  \item run the tests;
  \item fix any issues;
  \item commit your progress;
  \item push your work.
\end{itemize}

By the end of the session you will have written several C++ functions, explored defensive programming techniques, and used automated testing to check your work.

\begin{learningoutcomes}
By the end of this practical you should be able to:

\textbf{C++ Programming}
\begin{itemize}
  \item create and implement simple C++ functions;
  \item use function parameters and return values;
  \item distinguish between function declarations and function implementations;
  \item organise code using header \texttt{.h} and source \texttt{.cpp} files;
  \item use conditional statements to perform defensive checks on program inputs;
  \item read input from the user and display output to the console.
\end{itemize}

\textbf{Software Development Workflow}
\begin{itemize}
  \item clone and work with a repository provided through Classroom50;
  \item open and navigate an existing code base using Visual Studio Code;
  \item compile and run a C++ program;
  \item run automated tests using \texttt{make test};
  \item interpret test failures and use them to guide debugging;
  \item commit and push changes regularly using Git.
\end{itemize}

\textbf{Professional Practice}
\begin{itemize}
  \item understand the value of incremental development;
  \item recognise the importance of testing before committing code;
  \item appreciate the role of defensive programming in producing reliable software.
\end{itemize}
\end{learningoutcomes}

\begin{workflowreminder}
This workflow will be used throughout the module and will be particularly important when working on future practicals and assessments.

Whenever you are working on a programming task, aim to follow this cycle:

\begin{center}
\textbf{Read the task} $\rightarrow$ \textbf{Make a small change} $\rightarrow$ \textbf{Run the tests} $\rightarrow$ \textbf{Fix problems} $\rightarrow$ \textbf{Commit} $\rightarrow$ \textbf{Push} $\rightarrow$ \textbf{Repeat}
\end{center}

Useful commands:

\begin{bashcode}
git status
make test
git add .
git commit -m "Describe what you changed"
git push
\end{bashcode}

Try to commit regularly rather than waiting until the end of the practical. Frequent commits provide a useful record of your progress and help protect your work if something goes wrong.
\end{workflowreminder}

\CSCWorkflowDiagram


\begin{tipbox}
If the tests pass, that is a good sign. If the tests fail, they are trying to help you identify a problem. Read the error message carefully before making changes.
\end{tipbox}

\section{Getting Started with Classroom50}

In this practical, you will receive your work through Classroom50. This is the same workflow that will be used throughout the module and in future practical assessments.

\begin{importantnote}
\textbf{Windows Users}\par
\medskip

When this practical asks you to open a terminal, use the \textbf{CSC2059 Terminal} in Visual Studio Code.

On University laboratory computers, Visual Studio Code may open PowerShell by default. If this happens, open the terminal profile selector and choose:

\begin{terminal}
CSC2059 Terminal
\end{terminal}

Do not use PowerShell, Command Prompt, or Git Bash for the command-line activities in this practical.

On your own Windows computer, the CSC2059 Terminal should open automatically if you followed the \textbf{CSC2059 Development Environment: Windows Installation Guide}.
\end{importantnote}

\subsection{Step 1: Access your repository}

Open the link to the Lab Repo provided by your lecturer in Canvas.

If prompted, log in using your GitHub account. The system will automatically create a private repository for you containing the files required for this practical.

Once the repository has been created, open it and locate the \textbf{Code} button. Copy the repository URL.

\subsection{Step 2: Clone the repository}

You now need to download a local copy of your repository onto your computer.

\subsubsection{Option A: Using the command line}

Open a terminal and navigate to the folder where you would like to store your practical work. For example:

\begin{bashcode}
cd Documents
mkdir CSC2059
cd CSC2059
\end{bashcode}

Clone the repository:

\begin{bashcode}
git clone <YOUR_REPOSITORY_URL>
\end{bashcode}

Move into the repository directory.  It should look something similar to:

\begin{bashcode}
cd practical-02-<your GitHub username>
\end{bashcode}

Then open Visual Studio Code:

\begin{bashcode}
code .
\end{bashcode}

Visual Studio Code should open on the current folder.

\subsubsection{Option B: Using Visual Studio Code}

Open Visual Studio Code and select:

\begin{terminal}
View -> Command Palette
\end{terminal}

Search for:

\begin{terminal}
Git: Clone
\end{terminal}

Alternatively, choose \textbf{Clone Repository} from the welcome screen.

Paste your repository URL, choose a suitable folder on your computer, and when prompted select \textbf{Open Repository}. VS Code will open the practical automatically.


\subsection{Step 3: Explore the repository}

You should now see several files:

\begin{terminal}
README.md
Makefile
myFunctions.h
myFunctions.cpp
main.cpp
tests/public_tests.cpp
\end{terminal}

Take a moment to inspect these files. You do not need to understand everything immediately. For now, simply note that each file has a specific purpose.

\CSCRepositoryTree

\section{Repository Tour}

\subsection{Core files}

\begin{infobox}
\textbf{\texttt{README.md}}

The README file provides a brief overview of the practical, explains the repository structure, and contains useful commands and reminders.

Many software projects contain extensive documentation in their README file, so it is always worth taking a quick look before starting work.
\end{infobox}

\begin{infobox}
\textbf{\texttt{Makefile}}

The Makefile contains instructions that tell the computer how to compile and test the project.

Rather than typing long compiler commands every time, we can use shortcuts such as:

\begin{bashcode}
make run
make test
\end{bashcode}

You do not need to understand every detail of the Makefile yet. For now, think of it as a convenient shortcut that automates common tasks.
\end{infobox}

\begin{infobox}
\textbf{\texttt{myFunctions.h}}

This file contains function declarations, sometimes called prototypes.

A declaration tells the compiler that a function exists and describes:

\begin{itemize}
  \item its name;
  \item its parameters;
  \item its return type.
\end{itemize}

For example:

\begin{cppcode}
int iAdd(int a, int b);
\end{cppcode}

This tells us that the function is called \texttt{iAdd}, takes two integer parameters, and returns an integer.

Notice that declarations describe \textbf{what} a function looks like, but not \textbf{how} it works.
\end{infobox}

\begin{infobox}
\textbf{\texttt{myFunctions.cpp}}

This file contains the implementations of the functions declared in \texttt{myFunctions.h}. For example:

\begin{cppcode}
int iAdd(int a, int b)
{
    return a + b;
}
\end{cppcode}

The implementation contains the actual instructions that will be executed when the function is called.

As you work through this practical, most of your coding will take place in this file.
\end{infobox}

\begin{infobox}
\textbf{\texttt{main.cpp}}

This file contains a small C++ program that demonstrates how the functions in \texttt{myFunctions} can be used.

You can experiment with your functions here if you would like to try additional examples beyond those covered by the automated tests.

You can compile and run this program using:

\begin{bashcode}
make run
\end{bashcode}
\end{infobox}

\begin{infobox}
\textbf{\texttt{tests/public\_tests.cpp}}

This file contains the automated tests used throughout the practical.

Tests allow us to check whether our functions behave as expected. For example:

\begin{cppcode}
assert(iAdd(2, 3) == 5);
\end{cppcode}

This asks the computer to verify that the function produces the correct result. If the result is incorrect, the test will fail and provide information that can help us identify the problem.
\end{infobox}

\CSCCompilationModel

\subsection{Additional files}

Depending on the Classroom50 configuration, you may notice additional files or folders in your repository. Examples might include:

\begin{terminal}
.github/
.classroom/
workflows/
\end{terminal}

These files are used internally by Classroom50 and GitHub to support features such as automated testing, repository management, and assessment workflows.

\begin{importantnote}
For the purposes of this practical, you can safely ignore these files.

Unless explicitly instructed otherwise, you should not modify, rename, move, or delete any files that have not been discussed in this practical.

The files you will primarily work with are:

\begin{terminal}
myFunctions.h
myFunctions.cpp
main.cpp
tests/public_tests.cpp
\end{terminal}

If you are unsure whether a file should be edited, ask a demonstrator before making changes.
\end{importantnote}

\subsection{A quick observation}

\begin{importantnote}
You may notice that some functions already exist in \texttt{myFunctions.cpp}, but simply return placeholder values such as:

\begin{cppcode}
return 0;
\end{cppcode}

or:

\begin{cppcode}
return 0.0;
\end{cppcode}

This is intentional. The practical scaffold provides a starting point for your work.

Your task throughout the practical is to replace these placeholder implementations with working code, run the tests, and gradually make more tests pass.
\end{importantnote}

\begin{tipbox}
Do not be alarmed if some tests fail when you first run \texttt{make test}. At the beginning of the practical, this is entirely expected.
\end{tipbox}

\subsection{Step 4: Run the tests}

Open a terminal in VS Code:

\begin{terminal}
Terminal -> New Terminal
\end{terminal}

Run:

\begin{bashcode}
make test
\end{bashcode}

You may see something similar to:

\begin{terminal}
Running public tests...
Assertion failed:
iAdd(2,3) == 5
Aborted
\end{terminal}

This is expected. The tests are failing because some functions have not yet been implemented.

\begin{tipbox}
Think of the tests as a helpful assistant. They tell you what still needs attention.
\end{tipbox}

\section{Working Through the Exercises}

Each exercise in this practical follows the same process.

\subsection{Read the instructions}

Identify the function you are being asked to implement. For example:

\begin{cppcode}
int iAdd(int a, int b);
\end{cppcode}

\subsection{Implement the function}

Open:

\begin{terminal}
myFunctions.cpp
\end{terminal}

Locate the function. You might initially see something like:

\begin{cppcode}
int iAdd(int a, int b)
{
    return 0;
}
\end{cppcode}

Modify the implementation:

\begin{cppcode}
int iAdd(int a, int b)
{
    return a + b;
}
\end{cppcode}

Save the file.

\subsection{Run the tests}

Return to the terminal and run:

\begin{bashcode}
make test
\end{bashcode}

If the tests pass, you have successfully completed that part of the practical. If the tests still fail, read the error message carefully.

Compiler errors and test failures often contain useful information about what went wrong. A later section of this guide lists the most common errors you might encounter during the practical.

\begin{tipbox}
Learning how to search for, interpret and resolve programming errors is an important skill for any software developer. Even experienced developers spend a considerable amount of time reading documentation, consulting online resources, and investigating compiler messages.
\end{tipbox}

\subsection{Commit your progress}

Once you have completed an exercise and verified that the tests pass, commit your work.

\begin{bashcode}
git status
git add .
git commit -m "Implemented iAdd"
\end{bashcode}

Small, frequent commits are much easier to understand than a single large commit at the end of the practical.

\subsection{Continue}

Implement the next function, run the tests again, commit your progress, and repeat until all public tests pass.

\section{Submitting Your Work}

Classroom50 continuously tracks the contents of your repository. Strictly speaking, there is no separate \textbf{Submit} button.

However, your work only exists remotely if you have pushed it to GitHub. Before leaving the practical, make sure your latest changes have been uploaded.

Run:

\begin{bashcode}
git push
\end{bashcode}

You should see a message similar to:

\begin{terminal}
Everything up-to-date
\end{terminal}

or:

\begin{terminal}
Enumerating objects...
Writing objects...
Done.
\end{terminal}

As a final check, visit your repository in GitHub and verify that your latest commit is visible online.

\begin{importantnote}
Getting into the habit of pushing regularly is good practice and will be particularly important during practical assessments.
\end{importantnote}

\section{Common Errors and How to Interpret Them}

Programming can sometimes feel frustrating, particularly when you encounter error messages for the first time. The good news is that most errors are common, and with practice you will quickly learn how to recognise and fix them.

Compiler errors and test failures are useful sources of information. They are trying to tell you what went wrong.

\begin{reflectionbox}
If you encounter an error message, do not panic.

A useful habit is to ask yourself:

\begin{itemize}
  \item What is the error message telling me?
  \item Which file is involved?
  \item Which line number is mentioned?
  \item What changed since the last time the code worked?
  \item Have I run the tests again after making my changes?
\end{itemize}

Developing resilience when working with incomplete, incorrect, or failing code is an important skill for any software developer.
\end{reflectionbox}

\begin{commonerror}
\textbf{Error: expected \texttt{;}}

Example:

\begin{terminal}
error: expected ';' after expression
\end{terminal}

This usually means that a semicolon has been forgotten.

For example:

\begin{cppcode}
return a + b
\end{cppcode}

should be written as:

\begin{cppcode}
return a + b;
\end{cppcode}

Semicolons are used to mark the end of statements in C++.
\end{commonerror}

\begin{commonerror}
\textbf{Error: use of undeclared identifier}

Example:

\begin{terminal}
error: use of undeclared identifier 'iadd'
\end{terminal}

C++ is case-sensitive. These function names are considered different:

\begin{cppcode}
iAdd()
iadd()
IAdd()
\end{cppcode}

Check the spelling carefully and make sure the function name exactly matches the declaration in the header file.
\end{commonerror}

\begin{commonerror}
\textbf{Error: undefined symbols for architecture arm64}

Example:

\begin{terminal}
Undefined symbols for architecture arm64:
"divideXbyY(int, int)"
\end{terminal}

This means that the compiler has found the declaration of the function, but the linker cannot find its implementation.

Check:

\begin{itemize}
  \item Has the function been implemented?
  \item Is the function name spelled correctly?
  \item Do the parameter types match the declaration?
  \item Did you save the file before compiling?
\end{itemize}

For example:

\begin{cppcode}
int divideXbyY(int x, int y)
\end{cppcode}

is different from:

\begin{cppcode}
double divideXbyY(int x, int y)
\end{cppcode}
\end{commonerror}

\begin{commonerror}
\textbf{Error: assertion failed}

Example:

\begin{terminal}
Assertion failed:
iAdd(2,3) == 5
\end{terminal}

This means your program compiled successfully, but the behaviour of the function does not match what the test expects.

Read the error message carefully. The tests indicate which function still needs attention. Check the implementation and run the tests again.
\end{commonerror}

\begin{commonerror}
\textbf{Error: make: command not found}

Example:

\begin{terminal}
make: command not found
\end{terminal}

On Windows, first check that you are using the \textbf{CSC2059 Terminal} in Visual Studio Code.

Do not use PowerShell, Command Prompt, or Git Bash for CSC2059 command-line activities.

If you are working on your own Windows computer and have not yet configured the CSC2059 development environment, follow the separate \textbf{CSC2059 Development Environment: Windows Installation Guide}.

If this problem occurs on a University laboratory computer while using the CSC2059 Terminal, ask a demonstrator for assistance.
\end{commonerror}

\begin{commonerror}
\textbf{Error: fatal error: \texttt{myFunctions.h} file not found}

Example:

\begin{terminal}
fatal error: 'myFunctions.h' file not found
\end{terminal}

Check:

\begin{itemize}
  \item the spelling of the file name;
  \item capitalisation;
  \item the location of the file.
\end{itemize}

Remember that:

\begin{cppcode}
#include "myFunctions.h"
\end{cppcode}

is not necessarily the same as:

\begin{cppcode}
#include "myfunctions.h"
\end{cppcode}

on all operating systems.
\end{commonerror}

\begin{commonerror}
\textbf{Error: why are all my tests still failing?}

If every test still fails, ask yourself the following questions:

\begin{itemize}
  \item Did I save my changes?
  \item Did I edit the correct file?
  \item Am I working inside my repository folder?
  \item Did I implement the function I intended to implement?
  \item Did I run \texttt{make test} again after making changes?
\end{itemize}

Most issues can be solved by carefully reading the error message and working through the problem step by step.
\end{commonerror}

\section{Practical Exercises}

You will now complete the missing C++ implementations in the provided scaffold.

\begin{workflowreminder}
\textbf{Edit} $\rightarrow$ \textbf{Run \texttt{make test}} $\rightarrow$ \textbf{Read the result} $\rightarrow$ \textbf{Fix problems} $\rightarrow$ \textbf{Commit progress} $\rightarrow$ \textbf{Repeat}
\end{workflowreminder}

Most of your work will take place in:

\begin{terminal}
myFunctions.cpp
\end{terminal}

You may also edit:

\begin{terminal}
main.cpp
\end{terminal}

when asked to test interactive behaviour.

\exercise{Your First Function}{10 minutes}

Open:

\begin{terminal}
myFunctions.cpp
\end{terminal}

Find the function:

\begin{cppcode}
int iAdd(int a, int b)
{
    return 0;
}
\end{cppcode}

This function should return the sum of \texttt{a} and \texttt{b}. Update the function so that it behaves correctly.

Then run:

\begin{bashcode}
make test
\end{bashcode}

The first test should now pass.

Once you are happy with your work, commit your progress:

\begin{bashcode}
git add .
git commit -m "Implement iAdd"
git push
\end{bashcode}

\exercise{Defensive Division}{10 minutes}

The file contains two overloaded versions of \texttt{divideXbyY}:

\begin{cppcode}
double divideXbyY(double x, double y);
int divideXbyY(int x, int y);
\end{cppcode}

These functions should divide the first value by the second value.

However, division by zero is not allowed. Update both functions so that:

\begin{itemize}
  \item valid divisions return the correct result;
  \item division by zero returns \texttt{0} or \texttt{0.0}.
\end{itemize}

Run:

\begin{bashcode}
make test
\end{bashcode}

Commit your work once the relevant tests pass.

\exercise{Logarithms and Square Roots}{15 minutes}

Now implement:

\begin{cppcode}
double Log(double x);
double Sqrt(double x);
\end{cppcode}

You will need to use mathematical functions from the C++ standard library. The file already includes:

\begin{cppcode}
#include <cmath>
\end{cppcode}

The function \texttt{Log(x)} should return the natural logarithm of \texttt{x}. The function \texttt{Sqrt(x)} should return the square root of \texttt{x}.

However, both functions should be defensive:

\begin{itemize}
  \item \texttt{Log(x)} is only valid when \texttt{x > 0};
  \item \texttt{Sqrt(x)} is only valid when \texttt{x >= 0}.
\end{itemize}

If the input is invalid, return \texttt{0.0}.

Run:

\begin{bashcode}
make test
\end{bashcode}

Commit your work when the tests pass.

\exercise{A Simple Calculator Function}{20 minutes}

Implement:

\begin{cppcode}
double myCalculator(double a, char op, double b);
\end{cppcode}

This function should evaluate a simple expression of the form:

\begin{terminal}
a op b
\end{terminal}

where \texttt{a} and \texttt{b} are numbers, and \texttt{op} is one of \texttt{+}, \texttt{-}, \texttt{*}, or \texttt{/}.

For example:

\begin{cppcode}
myCalculator(3.0, '+', 2.0)
\end{cppcode}

should return:

\begin{terminal}
5.0
\end{terminal}

Your function should support:

\begin{terminal}
+
-
*
/
\end{terminal}

Remember to handle division by zero defensively. If the operator is not recognised, return \texttt{0.0}.

Run:

\begin{bashcode}
make test
\end{bashcode}

Commit your work when the calculator tests pass.

\exercise{Interactive Calculator}{15 minutes}

So far, the tests have called your functions automatically. Now you will call one of your functions from \texttt{main.cpp}.

Open:

\begin{terminal}
main.cpp
\end{terminal}

Modify the program so that it asks the user to enter a simple expression. For example:

\begin{terminal}
Enter a calculation: 3 + 2
\end{terminal}

The program should then call:

\begin{cppcode}
myCalculator(a, op, b)
\end{cppcode}

and print the result.

Example interaction:

\begin{terminal}
Enter a calculation: 3 + 2
Result: 5
\end{terminal}

You can run your program using:

\begin{bashcode}
make run
\end{bashcode}

\begin{importantnote}
This part is not checked by the public tests. You should test it manually using several inputs.
\end{importantnote}

Try:

\begin{terminal}
3 + 2
10 - 4
6 * 7
8 / 2
8 / 0
\end{terminal}

Commit your work once your interactive calculator behaves correctly.

\exercise{Sorting Three Numbers}{15 minutes}

Now implement:

\begin{cppcode}
void sortPrint(int i, int j, int k);
\end{cppcode}

This function should print the three numbers in descending order.

For example:

\begin{cppcode}
sortPrint(4, 9, 2);
\end{cppcode}

should print:

\begin{terminal}
9 4 2
\end{terminal}

You may use any approach you like. For example, you could use:

\begin{itemize}
  \item \texttt{if} statements;
  \item temporary variables;
  \item repeated comparisons.
\end{itemize}

This function prints directly to the console, so it is not currently checked by the public tests. To test it, add a call to \texttt{sortPrint} inside \texttt{main.cpp}, then run:

\begin{bashcode}
make run
\end{bashcode}

Try several examples:

\begin{cppcode}
sortPrint(1, 2, 3);
sortPrint(3, 2, 1);
sortPrint(2, 3, 1);
sortPrint(5, 5, 2);
\end{cppcode}

As an extension, modify the main program so that it asks the user to input three numbers to order. For example:

\begin{cppcode}
int i, j, k;
cout << "Enter three integers to order: ";
cin >> i >> j >> k;

sortPrint(i, j, k);
\end{cppcode}

Commit your work once you are satisfied.

\section{Optional Challenges}

If you complete the main exercises before the end of the practical, try one or more of the following challenges.

\begin{importantnote}
These challenges are not checked by the public tests. They are designed to give you additional practice.
\end{importantnote}

\begin{challengebox}
\textbf{Challenge A: Random Numbers and Loops}

Create a program in \texttt{main.cpp} that asks the user for two values:

\begin{terminal}
N R
\end{terminal}

where \texttt{N} is the number of random numbers to generate and \texttt{R} is the range of possible values.

The program should generate \texttt{N} random numbers in the range:

\begin{terminal}
0 to R - 1
\end{terminal}

and print their average.

You may find the following useful:

\begin{cppcode}
#include <cstdlib>
#include <ctime>
\end{cppcode}

and:

\begin{cppcode}
srand(time(0));
int rn = rand() % R;
\end{cppcode}

First write the program using a \texttt{for} loop. Then try rewriting it using:

\begin{itemize}
  \item a \texttt{while} loop;
  \item a \texttt{do...while} loop.
\end{itemize}

Remember that integer division truncates decimals. What happens if \texttt{sum} and \texttt{N} are both integers? How might you obtain a more accurate average?

Remember to commit your work.
\end{challengebox}

\begin{challengebox}
\textbf{Challenge B: Guessing Game}

Write a small guessing game in \texttt{main.cpp}. The program should:

\begin{enumerate}
  \item generate a random number between 0 and 9 inclusive;
  \item give the user three attempts to guess the number;
  \item print a suitable message if the user guesses correctly;
  \item print a suitable message if the user runs out of attempts.
\end{enumerate}

You may implement this using:

\begin{itemize}
  \item a \texttt{for} loop;
  \item a \texttt{while} loop;
  \item a \texttt{do...while} loop.
\end{itemize}

Try more than one version if you have time.

Have you committed your work?
\end{challengebox}

\begin{challengebox}
\textbf{Challenge C: Add Your Own Utility Function}

Add your own useful function to the \texttt{myFunctions} library.

Suggested examples:

\begin{cppcode}
int maximum(int a, int b);
int minimum(int a, int b);
bool isEven(int x);
double cube(double x);
int factorial(int n);
\end{cppcode}

Remember that adding a new function requires two steps.

First, add the declaration to:

\begin{terminal}
myFunctions.h
\end{terminal}

Then add the implementation to:

\begin{terminal}
myFunctions.cpp
\end{terminal}

Finally, call your function from \texttt{main.cpp} and test it using:

\begin{bashcode}
make run
\end{bashcode}

Remember to commit and push your work when finished.
\end{challengebox}

\section{Wrap-Up}

\begin{wrapupbox}
By now you have built a small C++ library from scratch, learned the header/source split that every C++ project uses, written your first defensive code, and pushed a submission through the same workflow you will use for the rest of this module.

Functions, files, and a bit of Git: that is a large part of the toolkit needed to begin building larger and more interesting programs.

Next time: putting these functions to work inside loops and data structures that are actually worth structuring.
\end{wrapupbox}


\end{document}
